\chapter{\ploren{Wprowadzenie}{Introduction}}
\label{chap:wprowadzenie}

% Zastąp poniższe wskazówki treścią właściwą dla pracy. Zachowaj logiczne
% przejście od szerszego kontekstu do problemu rozpatrywanego w pracy.

\begin{quote}
\itshape
\ploren
  {Przedstawiony układ wprowadzenia jest propozycją, a nie obowiązkowym schematem. W zależności od tematu, rodzaju pracy i ustaleń z promotorem poszczególne elementy można połączyć, rozwinąć albo ułożyć w innej kolejności, zachowując logiczny tok wypowiedzi. Tekst tej wskazówki należy usunąć z ostatecznej wersji pracy.}
  {The proposed structure of the introduction is a recommendation rather than a mandatory scheme. Depending on the subject, type of thesis and arrangements with the supervisor, individual elements may be combined, expanded or presented in a different order, provided that a clear logical progression is maintained. This guidance should be removed from the final version of the thesis.}
\end{quote}

\section{\ploren{Rola wprowadzenia}{Purpose of the introduction}}

\ploren
  {Wprowadzenie prowadzi od szerszego kontekstu dziedzinowego do konkretnego problemu technicznego, badawczego lub projektowego oraz wyjaśnia, dlaczego warto się nim zająć. Czytelnik powinien zrozumieć znaczenie tematu, zasadniczą trudność i logikę dalszej części dokumentu.}
  {The introduction proceeds from the broader disciplinary context to a specific technical, research or design problem and explains why it deserves attention. The reader should understand the significance of the subject, the principal difficulty and the logic of the remainder of the document.}

\ploren
  {Nie jest ono streszczeniem, rozbudowanym przeglądem literatury ani opisem wykonanych czynności. Szczegółowe cele, zakres, wymagania i kryteria oceny należą do kolejnego rozdziału, natomiast metodyka, wyniki i wnioski — do odpowiednich części pracy.}
  {It is not an abstract, an extensive literature review or a record of activities performed. Detailed objectives, scope, requirements and evaluation criteria belong in the following chapter, while the methodology, results and conclusions belong in their respective parts of the thesis.}

\section{\ploren{Kontekst i znaczenie tematu}{Context and significance of the subject}}

\ploren
  {Temat należy umieścić we właściwym obszarze elektrotechniki, automatyki, elektroniki, informatyki albo dyscypliny łączącej te dziedziny. Trzeba rozpocząć od informacji potrzebnych do zrozumienia pracy, a nie od ogólnej historii techniki. Kontekst może dotyczyć roli badanego urządzenia lub procesu, zmian technologicznych, wymagań przemysłowych albo konsekwencji określonego zjawiska.}
  {The subject should be placed within the relevant area of electrical engineering, automation, electronics, computer science or an interdisciplinary field. The discussion should begin with information needed to understand the thesis rather than with a general history of technology. The context may concern the role of the device or process under study, technological change, industrial requirements or the consequences of a particular phenomenon.}

\ploren
  {Oceny stanu techniki, dane i wymagania normatywne wymagają wiarygodnych źródeł. Określenia „ważny”, „nowoczesny” lub „dynamicznie rozwijający się” nie zastępują uzasadnienia opartego na konkretnych potrzebach, ograniczeniach albo korzyściach.}
  {Assessments of the state of technology, data and normative requirements require reliable sources. Descriptions such as “important”, “modern” or “rapidly developing” do not replace a justification based on specific needs, constraints or benefits.}

\section{\ploren{Problem i jego uwarunkowania}{Problem and its conditions}}

\ploren
  {Następnie należy wskazać, czego dotyczy problem, w jakich warunkach występuje, jakie powoduje konsekwencje oraz dlaczego dostępne rozwiązanie jest niewystarczające. Wystarcza zwięzłe sformułowanie; precyzyjna definicja, stan początkowy i oczekiwany rezultat należą do rozdziału dotyczącego celu i zakresu.}
  {The introduction should then explain what the problem concerns, under which conditions it occurs, what consequences it causes and why the available solution is insufficient. A concise statement is sufficient; the precise definition, initial situation and expected outcome belong in the chapter on the aim and scope.}

\ploren
  {Problem trzeba odróżnić od zadania. Zdanie „problemem jest wykonanie aplikacji” nie wskazuje trudności, którą aplikacja ma rozwiązać. Może nią być na przykład niedostateczna dokładność pomiaru, brak współpracy urządzeń, długi czas obliczeń, mała odporność algorytmu albo brak wiarygodnej walidacji.}
  {The problem must be distinguished from the task. The statement “the problem is to develop an application” does not identify the difficulty that the application is intended to solve. That difficulty may be insufficient measurement accuracy, lack of device interoperability, excessive computation time, limited algorithm robustness or insufficient validation.}

\section{\ploren{Motywacja podjęcia pracy}{Motivation for the thesis}}

\ploren
  {Motywacja wyjaśnia praktyczną, techniczną lub poznawczą wartość rozwiązania problemu. Może wynikać z potrzeby poprawy dokładności, niezawodności, bezpieczeństwa, efektywności energetycznej, wydajności, interoperacyjności albo kosztu. Można też krótko opisać rzeczywistą potrzebę przedsiębiorstwa lub laboratorium, nie ujawniając informacji poufnych.}
  {The motivation explains the practical, technical or scientific value of solving the problem. It may arise from a need to improve accuracy, reliability, safety, energy efficiency, performance, interoperability or cost. An actual need of a company or laboratory may also be described briefly without disclosing confidential information.}

\ploren
  {Osobiste zainteresowanie autora nie uzasadnia znaczenia tematu dla odbiorcy. Nie należy też przed badaniami deklarować przewagi projektowanego rozwiązania; można jedynie wskazać oczekiwaną korzyść, która zostanie później zweryfikowana.}
  {The author's personal interest does not demonstrate the significance of the subject to the reader. Nor should the superiority of the proposed solution be claimed before evaluation; only an expected benefit to be verified later may be stated.}

\section{\ploren{Krótkie osadzenie w stanie wiedzy}{Brief positioning within the state of knowledge}}

\ploren
  {Można zwięźle wskazać główne klasy istniejących rozwiązań i ograniczenie uzasadniające dalszą pracę, opierając opis na kilku bezpośrednio związanych źródłach. Nie należy tu kolejno streszczać publikacji ani przedstawiać szczegółowych podstaw teoretycznych.}
  {The main classes of existing solutions and the limitation that justifies further work may be identified briefly using a small number of directly relevant sources. Publications should not be summarised one after another here, nor should detailed theoretical background be presented.}

\ploren
  {Porównanie metod, podstawy teoretyczne i wyprowadzenie luki należą do rozdziału przeglądowego. Tutaj wystarcza logiczne przejście: istnieje potrzeba, dostępne podejścia mają istotne ograniczenie, dlatego zasadne jest podjęcie problemu. Luka może dotyczyć warunków badań, kosztu, braku implementacji, danych albo walidacji.}
  {The comparison of methods, theoretical foundations and derivation of the gap belong in the review chapter. Here, a logical transition is sufficient: a need exists, available approaches have a relevant limitation and therefore the problem warrants investigation. The gap may concern experimental conditions, cost, implementation, data or validation.}

\section{\ploren{Przedmiot pracy i zapowiedź dalszych rozdziałów}{Subject of the thesis and outline of subsequent chapters}}

\ploren
  {Przed zakończeniem warto wskazać przedmiot pracy: obiekt, proces, system, metodę albo zjawisko, którego dotyczy opracowanie. Nie należy jednak powtarzać celu głównego, celów szczegółowych, zakresu ani wymagań z następnego rozdziału.}
  {Before concluding, the subject of the thesis—the object, process, system, method or phenomenon considered—may be identified. The main objective, specific objectives, scope and requirements from the following chapter should not be repeated.}

\ploren
  {Ostatni akapit omawia rolę kolejnych rozdziałów, zamiast tylko powtarzać ich tytuły. Powinien pokazać przejście od przeglądu i wyboru podejścia przez metodykę oraz realizację do wyników i wniosków. Opis trzeba dostosować do faktycznego układu pracy.}
  {The final paragraph should explain the role of the subsequent chapters rather than merely repeat their titles. It should show the progression from the review and choice of approach through the methodology and implementation to the results and conclusions. The outline must reflect the actual structure of the thesis.}

\section{\ploren{Najczęstsze błędy}{Common mistakes}}

\ploren
  {Typowe błędy to: zbyt ogólne otwarcie, brak źródeł, utożsamianie problemu z zadaniem, powtarzanie celu i zakresu, przedwczesne wyniki, niezweryfikowane deklaracje oraz mechaniczne streszczanie spisu treści. Wprowadzenie warto ostatecznie zredagować po ukończeniu zasadniczej części pracy, aby odpowiadało faktycznej zawartości dokumentu.}
  {Common errors include an overly general opening, missing sources, confusing the problem with the task, repeating the aim and scope, premature results, unverified claims and mechanically summarising the table of contents. The introduction should be edited after the main body is complete so that it reflects the actual document.}
